% ============================================================================
% §1 INTRODUCCIÓN
% ============================================================================
\section{Introducción}
\label{sec:introduccion}

% ----------------------------------------------------------------------------
\subsection{Contexto y motivación}
\label{subsec:contexto}

Los mercados financieros constituyen uno de los sistemas dinámicos más estudiados y, a la vez, menos predecibles con los que trabaja la matemática aplicada. Cada sesión bursátil, un inversor debe decidir si comprar, vender o mantener un determinado activo, disponiendo únicamente de información histórica y de una estimación, necesariamente imperfecta, de la distribución futura de precios. Este problema, conocido en la literatura como \emph{trading cuantitativo}, puede formularse con rigor como un \emph{proceso de decisión secuencial bajo incertidumbre}: el agente no conoce el estado futuro del sistema, pero puede asignar probabilidades a sus posibles evoluciones y tomar decisiones que maximicen el valor esperado de un objetivo económico.

Desde esta perspectiva, el problema admite un tratamiento natural mediante herramientas que el estudiante de la asignatura \emph{Métodos Numéricos} ha encontrado en otros contextos: generación de variables aleatorias, simulación estocástica y estimación del valor esperado mediante promedios muestrales. El presente trabajo se apoya en los apuntes de Aràndiga~\cite{arandiga} sobre el método de Monte Carlo como soporte teórico y los extiende mediante el algoritmo \emph{Monte Carlo Tree Search} (MCTS), un método de planificación sobre árboles de decisión introducido por Kocsis y Szepesvári~\cite{kocsis2006} que combina muestreo estocástico con una política de exploración basada en cotas de confianza.

% ----------------------------------------------------------------------------
\subsection{Planteamiento del problema}
\label{subsec:problema}

Sea $\{P_t\}_{t=0}^{T}$ la serie temporal de precios de cierre diarios del activo TSLA (Tesla, Inc.) durante un horizonte de dos años. Un agente dispone de un capital inicial $C_0 = 10\,000$\,\$ y, al final de cada sesión $t$, debe ejecutar una acción $a_t$ perteneciente al conjunto finito
\[
    \mathcal{A} = \{\text{COMPRAR\_100},\ \text{COMPRAR\_75},\ \ldots,\ \text{MANTENER},\ \ldots,\ \text{VENDER\_100}\},
\]
que especifica la fracción del efectivo disponible destinada a comprar acciones, la fracción de acciones a liquidar, o la ausencia de operación. La decisión modifica el estado de la cartera $(\text{efectivo}_t,\ \text{acciones}_t)$, cuyo valor total se actualiza con el precio del día siguiente.

La pregunta operativa que guía el trabajo es la siguiente:

\begin{quote}
    \itshape ¿Existe una política $\pi: \mathcal{S} \to \mathcal{A}$, dependiente del histórico de precios y de la posición actual, que obtenga sistemáticamente un retorno superior al de la estrategia pasiva \emph{Buy \& Hold}?
\end{quote}

La estrategia \emph{Buy \& Hold} consiste en invertir todo el capital en el activo en $t=0$ y no realizar ninguna operación posterior; sirve como referencia neutra frente a la cual evaluar cualquier estrategia activa. La presente memoria no pretende responder a esta pregunta en toda su generalidad, sino construir y analizar un \emph{método concreto} ---MCTS con rollouts basados en movimiento Browniano geométrico--- y medir su comportamiento empírico sobre un periodo histórico específico.

% ----------------------------------------------------------------------------
\subsection{Por qué Monte Carlo Tree Search}
\label{subsec:por-que-mcts}

El espacio de estados de nuestro problema es continuo (el precio $P_t$ toma valores en $\mathbb{R}_{>0}$) y de dimensión apreciable, pues el estado no depende solo del precio actual sino también de indicadores técnicos como la media móvil y el RSI, y de la propia composición de la cartera. Esta combinación hace inviable una enumeración exhaustiva al estilo de la programación dinámica clásica. Alternativas como el \emph{Q-learning} tabular requieren una discretización costosa, mientras que las reglas heurísticas (cruces de medias, umbrales de RSI) carecen de garantías asintóticas.

MCTS ofrece un punto intermedio atractivo. En lugar de aproximar una función de valor sobre todo el espacio de estados, construye un árbol de búsqueda \emph{local} al estado actual y lo expande selectivamente mediante simulaciones aleatorias. La política de selección \emph{UCB1} (\emph{Upper Confidence Bound}) equilibra dos objetivos opuestos:
\begin{itemize}
    \item \textbf{Explotar} las acciones que han dado buen resultado en simulaciones anteriores.
    \item \textbf{Explorar} las acciones con pocas visitas, por si escondieran retornos altos aún no descubiertos.
\end{itemize}
Este compromiso, formalizado por la desigualdad de Chernoff--Hoeffding sobre las visitas y los retornos acumulados, garantiza que el algoritmo converge al \emph{minimax} del árbol de decisión a medida que el número de iteraciones crece~\cite{kocsis2006}.

% ----------------------------------------------------------------------------
\subsection{Conexión con el Monte Carlo clásico}
\label{subsec:puente-arandiga}

El aporte didáctico central de este trabajo reside en mostrar que \emph{MCTS no es un algoritmo esotérico del ámbito de la inteligencia artificial}, sino una aplicación directa del método de Monte Carlo por media muestral estudiado en el tema correspondiente de la asignatura. Si denotamos por $X_i$ el retorno económico acumulado al simular una trayectoria de precios durante $H$ días a partir del estado actual tras ejecutar una acción $a$, y se repite esta simulación $n$ veces de forma independiente, el estimador
\begin{equation}
    \widehat{I}_n \;=\; \frac{1}{n}\sum_{i=1}^{n} X_i \;\xrightarrow{\ \text{c.s.}\ }\; \mathbb{E}[X \mid a]
    \qquad (n \to \infty)
    \label{eq:mc-estimator-intro}
\end{equation}
converge casi seguramente al retorno esperado bajo la acción $a$, en virtud de la Ley Fuerte de los Grandes Números (Teorema~\ref{teorema:llng} de~\S\ref{sec:monte-carlo}).

El lector puede constatar que el cociente $Q(a)/N(a)$ que MCTS mantiene en cada nodo de su árbol no es otra cosa que el estimador~\eqref{eq:mc-estimator-intro} aplicado al conjunto particular de rollouts que el algoritmo haya ejecutado hasta el momento. La Sección~\ref{sec:monte-carlo} desarrollará esta identificación de forma rigurosa y derivará, como corolario, la cota de error $\mathcal{O}(n^{-1/2})$ que justifica la elección del parámetro \texttt{ITERACIONES}.

% ----------------------------------------------------------------------------
\subsection{Objetivos del trabajo}
\label{subsec:objetivos}

El presente documento persigue los cinco objetivos siguientes:

\begin{enumerate}
    \item[(O1)] \textbf{Formalizar} el problema de \emph{trading} diario como un proceso de decisión de Markov (MDP) con espacio de acciones discreto y recompensa financiera explícita.
    \item[(O2)] \textbf{Implementar} el algoritmo MCTS con política UCB1 y rollouts generados mediante un modelo de movimiento Browniano geométrico calibrado sobre ventana móvil.
    \item[(O3)] \textbf{Demostrar} que la convergencia del algoritmo se apoya directamente en la Ley Fuerte de los Grandes Números (Aràndiga,~\S3) y en la cota de error del Teorema del análisis del error (Aràndiga,~\S5).
    \item[(O4)] \textbf{Comparar} empíricamente la estrategia MCTS con la referencia \emph{Buy \& Hold} mediante métricas estandarizadas: retorno total, ratio de Sharpe, alfa relativa, \emph{drawdown} absoluto y \emph{drawdown} del alfa.
    \item[(O5)] \textbf{Validar} experimentalmente la cota $\mathcal{O}(n^{-1/2})$ observando la evolución del valor estimado en la raíz del árbol en función de \texttt{ITERACIONES}.
\end{enumerate}

% ----------------------------------------------------------------------------
\subsection{Estructura del informe}
\label{subsec:estructura}

La Sección~\ref{sec:variables-aleatorias} recoge los conceptos probabilísticos necesarios, siguiendo de cerca el tema~\S2 de los apuntes~\cite{arandiga}: variables aleatorias discretas y continuas, método de la transformación inversa y generación de muestras. La Sección~\ref{sec:monte-carlo} desarrolla íntegramente los métodos de Monte Carlo por éxito-fracaso y por media muestral (\S3 y~\S5 de~\cite{arandiga}), demostrando las leyes de los grandes números y el análisis del error, que constituyen la base teórica del MCTS. La Sección~\ref{sec:mdp-mcts} formaliza el MDP, introduce UCB1 y describe el ciclo de cuatro pasos del MCTS, así como el modelo GBM utilizado en los rollouts. La Sección~\ref{sec:implementacion} presenta la implementación en Python y discute las principales decisiones de diseño. La Sección~\ref{sec:resultados} recoge los resultados experimentales sobre TSLA con seis gráficas comentadas. Finalmente, la Sección~\ref{sec:discusion} contiene la discusión y las conclusiones.
